Against uniform FVI, compiled witness is earlier in 3 of 30 common successful target/repetition comparisons, later in 27, and tied in 0.
Against error-driven FVI, compiled witness is earlier in 15 of 30 common successful target/repetition comparisons, later in 15, and tied in 0.
These are observed complete-prefix wall times, not a hardware-general speed law. The comparison uses the prospective targets rather than retrospectively selecting a favorable breakpoint. The full positive-tolerance partition, including one-sided and joint nonattainment, is reported in the supplement and machine-readable frontier.
The error-driven rule returns uniform grids in all 210 recorded date models; all 60 adaptive checkpoints have exactly the same continuations and actors as their uniform-FVI counterparts. Thus its additional observed cost is pilot and allocation overhead, not evidence of a realized nonuniform-grid improvement. The nonuniform verification theorem is proved and tested, but an effective spatial-adaptation advantage is not established by this catalogue.
One prospective tolerance illustrates the distinction between common-resolution speed and target work. At horizon three, price four and target two, witness crosses at N=32 while uniform FVI first crosses at N=64. Median complete-prefix times are 5.498 and 54.146 seconds, respectively. This is one economic cell with three timing repetitions. At target one, witness attains all twelve services and each FVI implementation attains six; none attains one half or one quarter at its cap.
